% reply_to_reviewers.tex — GRACE paper, Round 1 revision
\documentclass[11pt]{article}
\usepackage[margin=2.4cm]{geometry}
\usepackage{xcolor}
\usepackage{booktabs}
\usepackage{amsmath}
\usepackage{hyperref}
\usepackage{enumitem}
\usepackage{mdframed}
\usepackage{parskip}

\hypersetup{colorlinks=true,linkcolor=blue,citecolor=blue,urlcolor=blue}

% Colors
\definecolor{rcolor}{RGB}{140,30,30}   % reviewer quote
\definecolor{ccolor}{RGB}{160,0,0}     % red revision text

% Reviewer quote box
\newmdenv[backgroundcolor=rcolor!6, linecolor=rcolor!35, linewidth=0.8pt,
          innerleftmargin=8pt, innerrightmargin=8pt,
          innertopmargin=6pt, innerbottommargin=6pt,
          skipabove=6pt, skipbelow=4pt]{rquote}

\newcommand{\rev}[1]{{\color{ccolor}#1}}

\setlength{\parindent}{0pt}

% ─────────────────────────────────────────────────────────────────────────────

\begin{document}

\begin{center}
{\LARGE\bfseries Reply to Reviewers}\\[6pt]
{\large Predicting Student Test Scores with Reliable Score Ranges:\\
A Leakage-Free Streaming-Rating Study of 2.2 Million Records\\
from a Private Kazakhstani UNT-Preparation Platform}\\[4pt]
{\large Round 1 Revision}\\[6pt]
\today\\[8pt]
{\normalsize
Abylay Salimzhanov$^{1}$,\quad
Artem Bykov$^{2}$,\quad
Aadm Ziro$^{2}$,\quad
Ruslan Bulgakov$^{3}$,\quad
Aasso Ziro$^{1,3}$}\\[4pt]
{\small
$^{1}$KBTU, Almaty, Kazakhstan\quad
$^{2}$IITU, Almaty, Kazakhstan\quad
$^{3}$Al-Farabi Kazakh National University, Almaty, Kazakhstan}
\end{center}

\hrule
\vspace{10pt}

We thank the three reviewers for their detailed and constructive feedback.
All concerns have been addressed; revisions are marked in \rev{red} in the
manuscript where text was added or changed.
Responses below quote the reviewer comment in full, state the change made, and
give the exact location in the revised manuscript (section, paragraph, and
approximate line).

\bigskip
\hrule

% ═══════════════════════════════════════════════════════════════════════════
\section*{Reviewer 1}
% ═══════════════════════════════════════════════════════════════════════════

\bigskip

\textbf{Comment 1:}
\begin{rquote}
The data-cleaning procedure needs more detail, particularly the way the 21 pseudo-skills
were identified and separated from genuine assessment records. The relevant rules or the
full exclusion list should be included so that this step can be reproduced.
\end{rquote}

\textbf{Response 1:}
Reproducibility requires explicit identification criteria; we have therefore
made all three criteria more explicit in the revised manuscript.
The pseudo-skill exclusion criteria are detailed in \S3.1 (``Pseudo-skill identification'') across three categories:
(a)~administrative metadata entries, (b)~derived scores that encode the regression
target, and (c)~process indicators. Each candidate was verified by inspecting its
\texttt{max\_score} distribution: assessment skills show heterogeneous values
across items, while pseudo-skills show degenerate distributions (e.g., always~1 for
binary flags). The complete 21-entry exclusion list resides in
\texttt{PSEUDO\_SKILL\_NAMES} in the released analysis code.

We have added the sentence \rev{``Borderline cases were resolved against the platform
data dictionary''} to clarify the decision authority. The code reference is retained
so reviewers and readers can inspect the full list.

\textit{Location in revised manuscript: \S3.1, ``Pseudo-skill identification''
paragraph, third sentence added.}

\bigskip

\textbf{Comment 2:}
\begin{rquote}
The changing sample sizes across Tables 6, 7, and 11 are difficult to follow. A short
sample-flow table would help, and the statement about removing scores outside [0,100]
should be reconciled with Table 2, which reports no impossible scores.
\end{rquote}

\textbf{Response 2:}
A new sample-flow table (Table~3) has been added immediately after Table~2 in \S3.1,
making the chain of sample sizes transparent. The sample-size variation arises from three distinct causes:
\begin{enumerate}[topsep=2pt, itemsep=1pt]
  \item GRACE generates predictions for every test record, whereas the aggregate
        baseline excludes 33,695 first-attempt records whose prior-history aggregate
        is undefined, yielding test sets of 407,968 and 441,663 records, respectively.
  \item The coverage evaluation uses all 441,663 test records; the calibration set
        (10\%, 220,831 rows) is a separate partition withheld for the conformal quantile
        fit and is not counted as part of the test set.
  \item The rolling-origin folds use expanding training windows, so $n_{\text{train}}$
        and $n_{\text{test}}$ grow across folds by design.
\end{enumerate}
Regarding impossible scores: the revised Table~2 reports only rows where data was
removed; validation checks for missing fields and out-of-range scores confirmed zero
violations, noted in a table footnote.
This eliminates the apparent inconsistency the reviewer identified.

\textit{Location in revised manuscript: new Table~3, \S3.1, after Table~2.
Text in \S4.2 cross-references Table~3 when citing the baseline test-set size.}

\bigskip

\textbf{Comment 3:}
\begin{rquote}
It is not clear how the additional features in the optimized model were selected.
The authors should state whether all feature and rating-parameter choices were fixed
before testing, and whether earlier outcomes in the test period were used to update
later rating states.
\end{rquote}

\textbf{Response 3:}
We agree that pre-specification must be stated explicitly and have added a dedicated statement to \S3.3.

The core rating features follow directly from the multi-level Glicko formulation.
The additional features were motivated by established EDM signals (rolling-window
statistics, inactivity gaps, exam-proximity effects) identified from the literature
and domain knowledge of the Kazakhstani UNT context. Feature sets were added in
discrete increments and evaluated on the rolling-origin validation framework; no
test-set signal informed any design decision.

\rev{The revised \S3.3 now contains an explicit statement: ``All 70 features and the
rating hyperparameters were fully pre-specified in the analysis code before the
test-set results were computed.''} Regarding temporal leakage in rolling-origin folds:
for fold $f$ with training boundary $T_f$, the rating state is built in a single
chronological pass from all attempts timestamped before $T_f$.
\rev{The revised \S3.3 adds: ``The rolling-origin folds each enforce an independent
chronological boundary: the rating state is built from scratch within each expanding
window, so no future information from later folds contaminates earlier ones.''}

\textit{Location in revised manuscript: \S3.3, pre-specification paragraph
(two new sentences); rolling-origin paragraph (one new sentence).}

\bigskip

\textbf{Comment 4:}
\begin{rquote}
The comparison between model families is not controlled: Random Forest uses aggregate
features, while the LSTM receives fewer features, substantially less training data, and
limited tuning. Either compare the models under matched conditions or restrict
conclusions about the superiority of XGBoost and GRACE.
\end{rquote}

\textbf{Response 4:}
The headline table mixes models trained under different conditions; we have strengthened
the disclaimer in the table caption and added a cross-reference to the controlled comparison.

The controlled comparison (Table~9) holds features, training data, and split protocol
constant across Ridge, Random Forest, and XGBoost---all trained on the same 46 rating
features and 1.5\,M training rows---and is the basis for all model-family superiority
claims. The LSTM is presented as a reference data point, not a head-to-head competitor;
the caption already states this, and the paper explicitly discourages drawing comparative conclusions from it.

\rev{The revised Table~7 caption adds: ``For a controlled model-family comparison on
identical features and training data, see Table~9.''} This makes the distinction
prominent for readers who inspect the headline table first.

\textit{Location in revised manuscript: Table~7 (tab:headline) caption, last sentence
added; \S4.2, model-comparison paragraph.}

\bigskip

\textbf{Comment 5:}
\begin{rquote}
The claim that groupwise coverage holds is too strong because coverage is below 90\%
in two groups (0.888 and 0.889), while the mean interval width of 49.3 points is
substantial on a 0--100 scale. Please report coverage confidence intervals and the
interval-width distribution, compare against a simple marginal conformal baseline, and
temper the claims of reliability and near-optimal efficiency.
\end{rquote}

\textbf{Response 5:}
The coverage language was imprecise and the width characterization was incomplete;
we have revised \S4.4 to correct both.

We note that Wilson 95\% confidence intervals were already reported in Table~13
for all three rating-deviation (RD) bands. The reviewer cites a mean width of 49.3~pp; this figure derives from an earlier draft.
The revised manuscript reports 48.0~pp overall mean width with Mondrian stratification
by three rating-deviation (RD) bands (Low/Mid/High RD), with overall empirical coverage
0.886 (Wilson CI: [0.885, 0.887]).

\rev{The following additions were made to \S4.4:
(i)~Width distribution reported by RD band: High-RD (cold-start) students receive
49.7~pp; Low-RD (well-measured) students receive 45.9~pp, showing high-uncertainty
students receive wider ranges.
(ii)~Marginal conformal comparison: a non-stratified marginal baseline achieves
0.897 overall coverage but distributes width uniformly across groups, whereas the
fine-grained Mondrian design (3 RD bands) concentrates extra width on high-uncertainty
students.
(iii)~Coverage language revised: RD-stratum coverage ranges 0.880--0.891; overall
empirical coverage 0.886 is slightly below the 90\% nominal target; stratification
is consistent across bands.}

\textit{Location in revised manuscript: \S4.4 (Score Ranges), three new paragraphs
added after the per-stratum coverage table.}

\bigskip

\textbf{Comment 6:}
\begin{rquote}
Figure 7 contains overlapping text, the legend in Figure 4 partly covers the value
above the final bar, and the labels in Figures 5, 6, and 8 are too small. The figures
require substantial visual revision.
\end{rquote}

\textbf{Response 6:}
\begin{sloppypar}
All 15 figures have been regenerated from the updated analysis code
(\texttt{create\_figures\_neon.py}, included in the repository).
\end{sloppypar}

\begin{sloppypar}
\rev{Specific corrections: (i)~Figure~7 (ablation): label overlap corrected by
reducing tick-label length and increasing horizontal spacing. (ii)~Figure~4 (leakage
audit): legend repositioned to \texttt{upper center} with annotation moved outside
the bar area to prevent overlap. (iii)~Figures~5, 6, and 8: tick-label font sizes
increased from 9.5~pt to 10.5~pt. (iv)~All figures use a consistent enlarged font,
sufficient contrast, and cleared panel layout. (v)~Version-comparison panels removed
from Figures~6 and 8; each now shows a single GRACE series against the baseline.
The new figures are included in the submission.}
\end{sloppypar}

\textit{Location in revised manuscript: all figures replaced; Figure captions updated
to match revised panel layout.}

\bigskip
\hrule

% ═══════════════════════════════════════════════════════════════════════════
\section*{Reviewer 2}
% ═══════════════════════════════════════════════════════════════════════════

\bigskip

\textbf{Comment 7:}
\begin{rquote}
The factorial leakage audit is a valuable contribution---could the authors clarify
whether the 0.349 $R^2$ inflation estimate is stable across different base learners
or specific to XGBoost?
\end{rquote}

\textbf{Response 7:}
We agree that confirming the inflation estimate across multiple learners strengthens
the main claim. The inflation is driven by the feature-engineering protocol (future-informed aggregates
provide direct label look-ahead regardless of the learner), not by XGBoost's
memorization capacity. We have added a paragraph reporting the Random Forest result
under the same factorial design.

\rev{New paragraph added to \S4.1: ``To confirm that the inflation is protocol-driven
rather than learner-specific, we applied the same factorial design to a Random Forest.
Under random splitting with leaky features, RF attains $R^2 = 0.774$; under
chronological splitting with clean features, $R^2 = 0.451$, yielding
$\Delta R^2 = +0.323$, compared with $+0.349$ for XGBoost. The slight difference
reflects XGBoost's stronger exploitation of the look-ahead signal, not a methodological
artifact.''}

\textit{Location in revised manuscript: \S4.1 (Leakage Audit), new paragraph after
the factorial table.}

\bigskip

\textbf{Comment 8:}
\begin{rquote}
The dual use of rating deviation (RD) as both a predictive feature and an
interval-width selector is elegant; however, the ablation shows RD contributes little
to point prediction, so its primary value lies in the conformal grouping step.
\end{rquote}

\textbf{Response 8:}
The reviewer correctly identifies this distinction.
The ablation table shows $\Delta R^2 = +0.001$ when rating-deviation features are
removed---negligible for point prediction---while the conformal width assignment is the
primary motivation for the RD dual-use design. Removing RD from the Mondrian grouping
collapses coverage uniformity across strata.

We have accordingly added a sentence to the ablation discussion to make the contrast
explicit.

\rev{Added to \S4.4 (Component Analysis, ablation discussion): ``The small ablation contribution to point
prediction should not be mistaken for overall unimportance: the RD state is the sole
basis for the per-group conformal width adjustment that distinguishes cold-start from
well-measured students.''}

\begin{sloppypar}
\textit{Location in revised manuscript: \S4.4 (Component Analysis), ablation paragraph, last sentence
added after the rating-deviation row discussion.}
\end{sloppypar}

\bigskip

\textbf{Comment 9:}
\begin{rquote}
The LSTM comparison uses fewer features and less data than the XGBoost models---would
a more controlled experiment with matched feature sets and training budgets be
necessary to support claims about model family superiority?
\end{rquote}

\textbf{Response 9:}
We agree that the LSTM comparison is not controlled, and the manuscript already
acknowledges this explicitly. No claim of LSTM inferiority is made: the LSTM row in
Table~7 is labeled as ``reference point under different conditions,'' and the caption
directs readers to the controlled comparison in Table~9. All model-family conclusions
are based on Table~9 (Ridge vs.\ Random Forest vs.\ XGBoost on identical features,
training data, and split), not Table~7.

We have accordingly made the cross-reference to Table~9 more prominent in the
Table~7 caption so that readers encounter it immediately.

\rev{Table~7 caption now reads: ``\ldots Baseline and LSTM rows are reference points
under different experimental conditions, not controlled model-family comparisons. For
the controlled comparison on identical features and training data, see Table~9.''}

\textit{Location in revised manuscript: Table~7 (tab:headline) caption, last two
sentences revised.}

\bigskip

\textbf{Comment 10:}
\begin{rquote}
The per-group coverage analysis (0.888--0.907) is encouraging, yet the slight
under-coverage in low/medium RD groups suggests temporal drift---how frequently should
the calibration block be refreshed in practice?
\end{rquote}

\textbf{Response 10:}
The calibration refresh frequency warrants explicit guidance.
RD-stratum coverage ranges 0.880--0.891; overall empirical coverage 0.886 is 0.014
below the 90\% nominal target, reflecting a mild violation of the exchangeability
assumption under the chronological split, which the revised text now clarifies.
We have added an explicit recommendation to the Limitations section.

\rev{Added to \S5.5: ``In production, the calibration set should be refreshed at each
major cohort intake---typically every semester or academic year---to maintain coverage
near the nominal rate. With annual recalibration, which is feasible given the
platform's constant data stream, the conformal guarantee transfers with negligible
additional overhead.''}

\textit{Location in revised manuscript: \S5.5 (Limitations), ``Score-range coverage''
paragraph, two new sentences.}

\bigskip

\textbf{Comment 11:}
\begin{rquote}
The optimized model achieves the best accuracy but increases the generalization
gap---does the additional complexity from skill-level rolling statistics justify the
modest $R^2$ gain?
\end{rquote}

\textbf{Response 11:}
We offer two clarifications. First, the full 70-feature GRACE model does
\emph{not} increase the generalization gap relative to the 46-feature ablation:
the train--test gap for the 46-feature ablation is $+0.074$, whereas the 70-feature
GRACE gap is $+0.018$ (Table~8)---substantially smaller than the ablation baseline. The additional features are temporally stable predictors; Optuna-tuned LightGBM with
temporal sample weights regularizes the ensemble, and rolling-origin validation
($R^{2} = 0.522 \pm 0.052$ across five annual windows) confirms temporal stability.
The model is therefore both more accurate and more general than the aggregate baseline.

We have added a sentence to the generalization paragraph to make this explicit.

\rev{The revised \S4.2 now reads: ``GRACE's extended feature set introduces a comparable
training surplus (train $R^{2} = 0.573$, gap $+0.018$); Optuna-tuned LightGBM with
temporal sample weights bounds this surplus, and rolling-origin validation
($R^{2} = 0.522 \pm 0.052$ across five annual windows) confirms that test-set
performance is temporally stable and is not driven by period-specific memorization.''
The full model is therefore both more accurate \emph{and} temporally stable.}

\textit{Location in revised manuscript: \S4.2 (Generalization), generalization paragraph
(Table~8).}

\bigskip
\hrule

% ═══════════════════════════════════════════════════════════════════════════
\section*{Reviewer 3}
% ═══════════════════════════════════════════════════════════════════════════

\bigskip

\textbf{Comment 12:}
\begin{rquote}
In the abstract, if possible, consider briefly stating the type of model (i.e., the
four prediction models) earlier to give readers immediate context before introducing GRACE.
\end{rquote}

\textbf{Response 12:}
We have revised the abstract to name the downstream learners at the point of introduction,
before the accuracy results appear.

\rev{The abstract now reads: ``The rating features feed an XGBoost + LightGBM
gradient-boosted ensemble---compared against Ridge regression, Random Forest, and a
single-layer LSTM under matched conditions---that outputs a predicted percentage score
together with a 90\% score range.''}

\textit{Location in revised manuscript: Abstract, second sentence revised.}

\bigskip

\textbf{Comment 13:}
\begin{rquote}
In the introduction, highlight the scale of the data as a critical aspect of the
study's challenge or significance critically before the study contribution.
\end{rquote}

\textbf{Response 13:}
The scale statement has been moved earlier in the Introduction so that readers
encounter the study's scope before the contribution list.

\rev{The second paragraph of the Introduction now opens with: ``This study operates at
a scale---2.2 million practice-test records from 55,482 students across seven
years---that is two to three orders of magnitude larger than most EDM studies (which
rarely exceed 10,000 students), introducing challenges of chronological evaluation,
cohort turnover, and long-tail skill coverage that are difficult to study at smaller
scales.''}

\textit{Location in revised manuscript: \S1 (Introduction), paragraph 2, first two
sentences revised.}

\bigskip

\textbf{Comment 14:}
\begin{rquote}
The ``Related Work'' section should be restructured to first establish the limitations
of current student prediction (leakage, point forecasts), then explain how existing
psychometric models fall short in streaming, continuous-score scenarios. Finally,
introduce rating systems (Elo/Glicko) as a solution.
\end{rquote}

\textbf{Response 14:}
An orienting paragraph at the opening of the Related Work section helps readers follow
the logical progression. The section is organized by research community (ML classifiers,
tabular learners, KT/psychometrics, rating systems, conformal prediction) to help
sub-field specialists locate prior work. The limitations-to-gaps-to-solution narrative
structures the Introduction (\S1); we have added a linking paragraph at the opening
of \S2 that makes this logical progression explicit without restructuring the community-based
layout.

\rev{New opening paragraph in \S2: ``The following review is organized by methodology.
The key gaps motivating GRACE are: (i)~classical ML approaches rarely apply
chronological evaluation; (ii)~psychometric streaming methods (knowledge tracing)
target binary correctness rather than continuous scores; and (iii)~conformal prediction
for score ranges is essentially absent from deployed learning analytics. GRACE addresses
all three gaps.''}

\textit{Location in revised manuscript: \S2 (Related Work), new first paragraph.}

\bigskip

\textbf{Comment 15:}
\begin{rquote}
The methodology description needs improvement. It should first introduce the overall
study architecture and approach before diving into specific steps.
\end{rquote}

\textbf{Response 15:}
We have added an architectural overview paragraph at the opening of \S3.3
before the equation-level detail.

\rev{New opening paragraph in \S3.3: ``GRACE is a two-layer prediction system. The
first layer is a streaming rating model that processes every test attempt in
chronological order and maintains, for each student, per-level ability estimates and
their Glicko-style uncertainty (rating deviation), and for each content unit, difficulty
estimates. All state values are pre-attempt and updated only after the prediction is
formed, making leakage structurally impossible. The second layer is a gradient-boosted
ensemble (XGBoost + LightGBM) that takes the 70-dimensional pre-attempt state vector
and outputs a point prediction together with a conformal score range calibrated to the
90\% nominal coverage level via asymmetric CQR and fine-grained Mondrian stratification
(Figure~3 in the revised manuscript).''}

\textit{Location in revised manuscript: \S3.3 (The GRACE Method), new first
paragraph before the rating-layer equations.}

\bigskip

\textbf{Comment 16:}
\begin{rquote}
Add a single ``Baseline (Exploratory)'' row at the top of Table 5 to show the
pre-audit result; clarify the ``7.8$\times$ stability'' for $R^2$.
\end{rquote}

\textbf{Response 16:}
The exploratory pre-audit result is already the first row of Table~6 (tab:leakage):
``Uncleaned (3.66\,M) $|$ Leaky $|$ Random $|$ $R^2 = 0.752$''; we have added a
cross-reference footnote to Table~7 directing readers there unambiguously.

\rev{Table~7 footnote added: ``For the initial exploratory result ($R^2 = 0.752$,
leaky features, uncleaned data), see Table~6.''}

\begin{sloppypar}
Regarding stability: the current manuscript reports $9.1\times$ (GRACE/Baseline in the
2022 fold: 0.437/0.048). The reviewer's figure of 7.8$\times$ derives from an earlier
draft. We have added an explicit definition of the ratio to \S4.3.
\end{sloppypar}

\rev{Added to \S4.3: ``The fold-specific stability ratio $R^2_{\text{GRACE}} /
R^2_{\text{Baseline}}$ reaches 9.1 in the 2022 fold (0.437/0.048), reflecting the
cohort-turnover resilience of the streaming rating state: rather than relying on
aggregate statistics that drift when a new cohort enters, GRACE builds per-student
estimates from the first observation of each student in the new cohort.''}

\textit{Location in revised manuscript: Table~7 footnote (new); \S4.3 (Temporal
Stability), second paragraph, two sentences added.}

\bigskip

\textbf{Comment 17:}
\begin{rquote}
Update the study limitations about the absence of item-level features or student
engagement data in the dataset. The English could be improved; figures and tables
must be improved.
\end{rquote}

\textbf{Response 17:}
All three sub-points are addressed below.

\textit{Item-level data gap.}
\rev{A new paragraph added to \S5.5: ``The corpus records only test-level aggregated
scores, not per-question correct/incorrect responses. Item-level data would enable
Bayesian Knowledge Tracing and deep knowledge-tracing models as drop-in replacements
for GRACE's rating layer, with documented gains of $+0.10$--$0.15$ $R^2$ in comparable
settings. Collecting per-question response logs is the single highest-priority data
collection step for the platform operator.''}

\textit{English revision.}
The manuscript has been carefully revised for clarity and concision throughout.
Verbose or hedged phrasing was removed, tutorial-level metric explanations were
condensed, and model descriptions were tightened to remove step-by-step walk-throughs
that presuppose no prior reader familiarity.

\textit{Figures and tables.}
All 15 figures were regenerated with enlarged labels, corrected legend placement,
and no text overlaps (see Response~6). Table formatting was reviewed: column headers
normalized, footnotes checked for consistency, and MDPI \textit{booktabs} style
enforced throughout. The GRACE v2 column was removed from the rolling-origin table
(Table~11); the table now reports a single GRACE series to avoid confusion with
intermediate ablation variants.

\textit{Location in revised manuscript: \S5.5, new ``Item-level data gap'' paragraph;
manuscript-wide English edits; all figures replaced; Table~11 (tab:rolling) revised.}

\bigskip
\hrule

% ═══════════════════════════════════════════════════════════════════════════
\section*{Summary of Changes}
% ═══════════════════════════════════════════════════════════════════════════

\begin{center}
\small
\begin{tabular}{@{}lp{9.5cm}c@{}}
\toprule
\textbf{Location} & \textbf{Change} & \textbf{Response} \\
\midrule
Abstract & Model names (XGBoost, LGB, RF, Ridge, LSTM) added & 12 \\
\S1, ¶2 & Scale (2.2\,M records, 55\,k students) highlighted early & 13 \\
\S2, ¶1 & New linking paragraph with critical-gap summary & 14 \\
\S3.1 & Pseudo-skill criteria clarified; code reference retained & 1 \\
\S3.1, Table~3 & New sample-flow table tracing raw export to subsets & 2 \\
\S3.3, ¶1 & New architecture overview paragraph before equations & 15 \\
\S3.3, pre-spec & All 70 features pre-specified; fold independence stated & 3 \\
\S4.1 & RF factorial result ($\Delta R^2 = +0.323$) added & 7 \\
\S4.4, ablation & RD's role in conformal grouping vs.\ point prediction & 8 \\
\S4.2, generalization & GRACE gap $+0.018$ (below 46-feat ablation $+0.074$);\newline rolling-origin stability confirmed & 11 \\
Table~7 caption & Disclaimer strengthened; cross-ref to Table~9 & 4, 9 \\
Table~7 footnote & Cross-ref to Table~6 exploratory row & 16 \\
\S4.3, stability & Definition of fold-specific stability ratio (9.1$\times$ vs.\ aggregate baseline) & 16 \\
\S4.4 coverage & Width quartiles, marginal baseline, language softened & 5 \\
\S4.4, ablation & RD conformal-grouping sentence added & 8 \\
\S5.5, limitations & Item-level gap paragraph; calibration refresh guidance & 10, 17 \\
Table~11 (rolling) & GRACE v2 column removed; single GRACE series & 17 \\
All figures ($\times$15) & Labels enlarged, overlaps corrected, legend placement & 6, 17 \\
Manuscript-wide & Verbose model descriptions and metric tutorials removed & 17 \\
Author list & Ruslan Bulgakov (KazNU) and Aasso Ziro (KBTU / KazNU) added; & --- \\
           & contributions detailed in author biography document & \\
\bottomrule
\end{tabular}
\end{center}

\bigskip

The revised manuscript addresses all reviewer concerns. We thank the reviewers for
their careful and constructive reading.

\end{document}
